\sect{客观题答案}
\begin{center}
\begin{tabular}{c*{10}{c}}
题号 & 1 & 2 & 3 & 4 & 5 & 6 & 7 & 8 & 9 & 10\\
\hline
答案 & B & D & A & C & C & B & A & D & B & D\\[3pt]
题号 & 11 & 12 & 13 & 14 & 15 & 16 & 17 & 18 & 19 & 20\\
\hline
答案 & 对 & 错 & 错 & 对 & 对 & 错 & 错 & 对 & 对 & 错\\
\end{tabular}
\end{center}

\sect{一、单项选择题详解}
\begin{enumerate}[label=\arabic*.]
\item \textbf{B。} 训练均值为 $5$，总体方差为 $(9+1+1+9)/4=5$，因此 z-score 为 $(10-5)/\sqrt5=\sqrt5$。min-max 值为 $(10-2)/(8-2)=4/3$。A 将 min-max 结果截到 $1$。C 用了样本方差 $20/3$。D 把测试值也计入统计量，违背题设。
\item \textbf{D。} 甲的 MAE、MSE 分别为 $2,12$，乙分别为 $3,9$。A 在 MSE 上选错，平方放大了甲的单个大误差。B 两项都选反。C 在 MAE 上选错，甲的绝对误差总和只有 $6$。
\item \textbf{A。} ID3、C4.5 和分类 CART 的常见划分准则分别如 A 项所述。依题设，错分数从 $3$ 降到 $2$，应剪掉该子树。B 对调了 ID3 和 C4.5 的准则。C 把熵或基尼替换为训练错分率。D 用训练子树的变化取代了指定的验证规则。
\item \textbf{C。} $P(X=x)=0.20$，三个后验依次为 $0.15,0.45,0.40$，最大的是类别 $2$。误判概率是其余两类后验之和 $0.55$。A 把最大后验误作错误概率。B 没选后验最大的类。D 的 $0.11$ 是在给定特征值与误判同时成立时的联合概率。
\item \textbf{C。} 在通常的独立同分布等大样本条件下，1-NN 的极限风险夹在 $R^\ast$ 和 $2R^\ast$ 之间。中心处的 $\ell_\infty$ 邻域体积为 $(2r)^d$，均匀分布时这也是邻域的概率质量。A 误将渐近结论用于每一个有限训练集。B 把风险上界改成相等，并忽略概率质量随维数减小。D 错把 $2R^\ast$ 当成下界，且说反了维数效应。
\item \textbf{B。} AdaGrad 的累计量为 $\sum_{s\le t}g_s^2$，RMSProp 的递推为 $v_t=\rho v_{t-1}+(1-\rho)g_t^2$，其中 $0<\rho<1$。早期梯度在前一个累计量中权重不变，在后一个累计量中权重衰减。A 把平方和误作带符号的一阶和。C 忽略了两种算法都使用历史梯度。D 把 RMSProp 的指数权重误作等权。
\item \textbf{A。} VGG 通过堆叠小卷积核来扩展网络，GoogLeNet 使用多分支 Inception，ResNet 用捷径学习残差。B 对调 VGG 和 GoogLeNet。C 将残差连接和小卷积堆叠的代表架构互换。D 对调了多分支与残差架构。
\item \textbf{D。} 训练误差本身偏大，训练与验证差距又小，支持高偏差的诊断。减小过大的 $k$ 可降低平滑程度，效果仍需验证。A 两条曲线接近不能排除欠拟合，题中已有可达的更低误差。B 把高偏差当成高方差。C 增加样本无法保证消除模型偏差。
\item \textbf{B。} 三类召回率为 $0.9,0.4,0.2$，宏平均为 $(0.9+0.4+0.2)/3=0.50$，微平均为 $(72+6+1)/(80+15+5)=0.79$。A 在宏平均中也按样本总数汇总。C 对调二者。D 在微平均中也按类别等权。
\item \textbf{D。} 给定 $X=0$ 或 $1$ 后，输出的熵分别为 $0,1$，故 $H(Y\mid X)=0.5$。$P(Y=1)=1/4$，所以 $H(Y)\approx0.8113$，$I(X;Y)\approx0.3113$。再用 $H(X)=1$ 得 $H(X\mid Y)=1-I(X;Y)\approx0.6887$。A 交换两个条件熵。B 将两个条件熵都替换成边缘熵。C 将互信息误算为两个边缘熵之差。
\end{enumerate}

\sect{二、判断题详解}
\begin{enumerate}[label=\arabic*.,start=11]
\item \textbf{对。} 每次测试的分割线形如 $x_1=t$ 或 $x_2=t$。它在祖先分出的区域内截出的边仍与坐标轴平行，决策边界由这些边的一部分组成。
\item \textbf{错。} 同一患者的影像可能共享解剖或成像特征。若两组都见过同一患者，就可能高估模型对新患者的表现。应以患者为划分单位。
\item \textbf{错。} 得分 $\bm u^\top\bm x$ 确实变号，但重构为 $\bm u\bm u^\top\bm x$。将两个 $\bm u$ 都取反后，线性投影矩阵保持原样。
\item \textbf{对。} $1\times1$ 卷积每个空间位置作一次通道间线性组合，能重新混合通道，但该层不读取邻近位置。
\item \textbf{对。} 对高斯核 $h^{-1}K((x-x_i)/h)$ 作换元，积分为 $\int K(u)\,\mathrm du=1$。取 $n$ 个这样的核的平均，积分仍为 $1$。
\item \textbf{错。} 前者是仅用当前及更早观测的滤波后验，后者还结合未来观测。后向量以 $i_t$ 为条件，通常因状态而异，乘进前向量后会改变归一化的后验。
\item \textbf{错。} $\varepsilon$ 增大时，每个点的邻域只能扩大或保持不变，原来达到 $\mathrm{MinPts}$ 的点仍达到该阈值，因此不会丢掉核心点身份。
\item \textbf{对。} 将原目标 $J=S+\lambda\|\bm w\|_1$ 整体乘 $1/n$，得到 $S/n+(\lambda/n)\|\bm w\|_1$。仅缩放数据项会改变正则项的相对权重。
\item \textbf{对。} 若行置换矩阵为 $P$，则 $Q,K,V$ 变为 $PQ,PK,PV$。得分矩阵按 $S\mapsto PSP^\top$ 重排，逐行 softmax 也按相同方式重排。再乘 $PV$，只留下输出行的置换 $P$。
\item \textbf{错。} 两组分别含 $2$ 和 $6$ 个样本。合并后的平均梯度应是两组平均值的 $1/4$ 与 $3/4$ 加权和，只有在两组平均梯度相同等特殊情形下，等权才等价。
\end{enumerate}

\sect{三、计算与分析题参考答案与评分标准}

\noindent 等价的公式和解法给同等分数。数值合理舍入不扣分。前问数值错误且后问沿用时，后问按其步骤给分。每问扣分不超过该问分值。

\noindent\textbf{1. 小批量分类网络的反向传播与归一化}\pts{10}

\begin{enumerate}[label=(\arabic*)]
\item 两个激活前输入依次为 $\ln 3,-\ln 3$，故 $h_A=3/4$，$h_B=1/4$。
两个输出各 $0.5$ 分。将 $\ln 3$ 与三类输出的 $r$ 混用，扣受影响结果的分值。
\item $\bm{z}_A=(r,0,-r)^{\mathsf T}$，$\bm{z}_B=(-r,0,r)^{\mathsf T}$，故
\[
\bm{p}_A=(4,2,1)^{\mathsf T}/7,\qquad \bm{p}_B=(1,2,4)^{\mathsf T}/7,
\qquad J=\ln 7-\tfrac12\ln 2\approx1.599.
\]
概率和损失各 $0.5$ 分。损失求和但未求平均，扣损失部分 $0.5$ 分。
\item 对单样本，得分梯度为概率向量减去该标签的独热向量，得到
$\bm{\delta}_A=(4,-5,1)^{\mathsf T}/7$ 和 $\bm{\delta}_B=(-6,2,4)^{\mathsf T}/7$。
两向量各 $0.5$ 分。减错独热向量位置或在此处提前除以 $2$，扣受影响的对应部分。
\item 以两条样本梯度的均值作为批次梯度，得到
\begin{align*}
\nabla_{\bm{v}}J&=(h_A\bm{\delta}_A+h_B\bm{\delta}_B)/2
=(6,-13,7)^{\mathsf T}/56,\\
\nabla_{\bm{c}}J&=(\bm{\delta}_A+\bm{\delta}_B)/2=(-2,-3,5)^{\mathsf T}/14.
\end{align*}
两组各 $0.5$ 分。偏置梯度多乘 $h$，扣偏置部分 $0.5$ 分。
\item 隐层误差信号为 $(\bm{v}^{\mathsf T}\bm{\delta})h(1-h)$。两个样本均有 $h(1-h)=3/16$，故 $\delta_{h,A}=9r/28$、$\delta_{h,B}=-15r/14$。
各 $0.5$ 分。漏乘 Sigmoid 导数，扣对应结果分值。
\item
\begin{align*}
\nabla_{\bm{w}}J&=(\delta_{h,A}\bm{x}_A+\delta_{h,B}\bm{x}_B)/2
=(9r/56,-15r/28)^{\mathsf T},\\
\partial J/\partial b&=(\delta_{h,A}+\delta_{h,B})/2=-3r/8.
\end{align*}
两组各 $0.5$ 分。把 $\delta_{h,A}$ 与 $\delta_{h,B}$ 之差作为偏置梯度，扣 $0.5$ 分。
\item $u_{b,\rm new}=r/8-3r/8=-r/4$，其他新动量等于各自的批次梯度，故
\begin{align*}
\bm{w}'&=\bm{w}-(9r/56,-15r/28)^{\mathsf T}/10,& b'&=r/40,\\
\bm{v}'&=\bm{v}-(6,-13,7)^{\mathsf T}/560,&
\bm{c}'&=\bm{c}-(-2,-3,5)^{\mathsf T}/140.
\end{align*}
隐藏层与输出层各 $0.5$ 分。漏加 $b$ 的旧动量，扣隐藏层部分 $0.5$ 分。
\item BatchNorm 对每列用两个样本统计，列均值为 $(2,5)$，列方差为 $(1,4)$。标准化结果为 $\begin{pmatrix}-1&-1\\1&1\end{pmatrix}$。
正确均值和方差 $0.5$ 分，结果 $0.5$ 分。按行统计，扣对应两部分。
\item LayerNorm 在每行内跨特征统计，两行均值依次为 $2,5$，方差为 $1,4$，结果为 $\begin{pmatrix}-1&1\\-1&1\end{pmatrix}$。这里两组统计数值碰巧相同，但 BatchNorm 沿批次方向，LayerNorm 沿同一样本的特征方向。
正确按行统计 $0.5$ 分，结果和轴的区别 $0.5$ 分。照搬上一问的输出，扣结果部分 $0.5$ 分。
\item 单样本 SGD 在处理完 $A$ 后立即改变参数，$B$ 的梯度在更新后的参数处计算。小批量把两条梯度都在同一组初始参数处计算再取均值。由于第二次求导的位置一般不同，两者的结果一般不同。
明确 SGD 的第二条梯度使用新参数得 $0.5$ 分，说明小批量在同一点取均值得 $0.5$ 分。仅写“两者批次大小不同”，不给分。
\end{enumerate}

\noindent\textbf{2. 支持向量机、异或与多项式核}\pts{8}
\begin{enumerate}[label=(\arabic*)]
\item 正类两点和负类两点的中点都是 $(0,0)^\top$。若某个仿射函数能严格分开两类，它在正类中点必须为正、在负类中点必须为负，产生矛盾。只写不可分而未说明理由扣 0.5 分。
\item $h_1=H(r_1+r_2-1/2)$、$h_2=H(r_1+r_2-3/2)$，分别检测至少一个、两个都开启（0.5 分）。输出 $H(h_1-h_2-1/2)$（0.5 分）。隐藏逻辑或输出逻辑错误，各扣对应的 0.5 分。等效网络同样给分。
\item 一种映射是
$\bm\varphi(\bm x)=(1,\sqrt{2}x_1,\sqrt{2}x_2,x_1^2,\sqrt{2}x_1x_2,x_2^2)^\top$。
展开得到常数项与 $2\bm x^\top\bm z$（0.5 分），平方项含两个平方积和 $2x_1x_2z_1z_2$（0.5 分）。漏常数或线性项扣前 0.5 分，交叉项或平方项系数错扣后 0.5 分。
\item $G=8I+\bm1\bm1^\top$，即对角为 9（0.5 分），其余均为 1（0.5 分）。对角或非对角项错误，各扣对应的 0.5 分。
\item
\[
\max_{\bm\alpha\ge0,\ \sum_i\alpha_iy_i=0}
\left(\sum_i\alpha_i-\frac12\sum_{i,j}\alpha_i\alpha_jy_iy_jG_{ij}\right).
\]
核对偶式含目标和约束得 0.5 分。$\bm1\bm1^\top$ 对二次项的贡献为 $(\sum_i\alpha_i y_i)^2=0$，剩下 $\sum_i(\alpha_i-4\alpha_i^2)$（0.5 分）。漏掉标签或等式约束扣前 0.5 分，化简系数错误扣后 0.5 分。
\item 每个严格凹的二次项单独在 $\alpha_i=1/8$ 处最大（0.5 分）。这组值同时满足非负和等式约束，故也是唯一约束最优解（0.5 分）。乘子错误扣前 0.5 分，只用对称性断言最优扣后 0.5 分。
\item 四点乘子均正，任选一点都有 $y_if(\bm x_i)=1$。取第一点，
$\sum_j\alpha_jy_jG_{j1}=-1$，故 $b=0$（0.5 分）。展开
$f(\bm z)=\frac18\sum_jy_j(1+\bm x_j^\top\bm z)^2=-z_1z_2$（0.5 分）。
偏置错误或漏用等号条件扣前 0.5 分，判别函数错误扣后 0.5 分。
\item $f(2,-1/2)=1$，判为 $+1$。分数与标签各 0.25 分，错则扣对应分数。
\item $C_{\min}=1/8$（0.5 分）。$C\ge1/8$ 时新可行域包含原最优点，缩小可行域后它仍最优。$C<1/8$ 时原解违反盒约束（0.5 分）。界值错扣前 0.5 分，没有解释两侧的可行性扣后 0.5 分。
\end{enumerate}

\noindent\textbf{3. 混合模型的更新与聚类评价}\pts{12}
\begin{enumerate}[label=(\arabic*)]
\item $\widehat p=k/(n\,2r)$，所以比值是 $2:1$。高斯窗过窄会使估计在样本附近形成尖峰、对抽样波动敏感，方差较大。
\textbf{评分} 密度比 0.5 分，带宽过小的表现 0.5 分。用面积代替一维区间长度扣 0.5 分。
\item 累计平均为 $\bigl[4(-2,0)+(1,1)+(1,-1)\bigr]/6=(-1,0)$。设 $v$ 的距离为 $D$，则 $D^2/(4+D^2)=9/13$，所以 $D=3$。
\textbf{评分} 两结果各 0.5 分。只用新批次均值、或把距离当成平方距离权重，各扣对应 0.5 分。
\item 设样本为 $(x,y)$。初始混合权重比为 $1/4$，两均值的平方距离差为 $(x+1)^2-(x-1)^2=4x$。因此
\[
 r_{i2}/r_{i1}=\tfrac14\exp(2x/\sigma^2)=2^{x-2}.
\]
按题目顺序，$r_{i2}$ 为 $1/9,1/9,1/3,1/3$。
\textbf{评分} 比值 0.5 分，四个责任度 0.5 分。漏先验比扣 0.5 分，归一化方法正确可保留责任度对应的方法分。
\item $N_1=28/9,\ N_2=8/9$。用 $\pi_k'=N_k/4$ 得 $(7/9,2/9)$。用 $\mu_k'=\sum_i r_{ik}x_i/N_k$ 得 $(-1/7,0)$、$(1/2,0)$。
\textbf{评分} 权重 0.5 分，均值 0.5 分。将责任度作硬分配的，扣受影响项的分。
\item 共享球形方差的分母为 $4\cdot 2=8$。
\[
 \sigma_{\mathrm{new}}^2
 =\frac{\sum_{i,k}r_{ik}\|x_i-\mu_k'\|^2}{8}
 =\frac{8-(28/9)(1/7)^2-(8/9)(1/2)^2}{8}
 =\frac{27}{28}.
\]
\textbf{评分} 更新式 0.5 分，结果 0.5 分。分母漏维数扣 0.5 分，后续若代入自洽，可给结果项的方法分。
\item 纯度 $(3+1+2)/8=3/4$。28 个无序对中，真正为同类的有 12 个，预测同簇的有 8 个，二者同时成立的有 4 个，二者均不成立的有 12 个。因此兰德指数 $(4+12)/28=4/7$。
\textbf{评分} 每个指标 0.5 分。仅用同类且同簇的对计算兰德指数扣 0.5 分。
\item $TP=4,\ FP=4,\ FN=8$，故精确率 $1/2$，召回率 $1/3$，$F_1=2/5$。
\textbf{评分} 正确区分精确率与召回率 0.5 分，$F_1$ 0.5 分。调换两率扣 0.5 分。
\item 样本比例按表归一化。用 $I(Y;C)=\sum_{y,c}p_{yc}\log_2[p_{yc}/(p_y p_c)]$，非零贡献为 $(3/8)\log_2(3/2)$、$(1/8)\log_2(1/2)$ 和 $(1/4)\log_2 2$，故 $I=(3/8)\log_2(3/2)+1/8\approx0.3444$。
\textbf{评分} 互信息计算式 0.5 分，代入与结果 0.5 分。以表内计数直接当概率扣 0.5 分。
\item $H(Y)=1,\ H(C)=3/2$。同质性 $I/H(Y)=I\approx0.3444$，完整性 $I/H(C)=2I/3\approx0.2296$，两者的调和平均 $V=4I/5\approx0.2755$。
\textbf{评分} 两个归一化指标 0.5 分，调和平均 0.5 分。颠倒分母或用算术平均各扣对应 0.5 分。
\item 单链取最小值 3，全链取最大值 10，平均链取四个跨簇点对的平均 $27/4$。
\textbf{评分} 单链、全链各 0.25 分，平均链 0.5 分。将子簇 $A$、$B$ 的平均距离等权组合扣 0.5 分。
\item 甲是均值偏移。样本从各自初值沿更新走到密度峰附近，收敛到同一模式的样本可归为同一簇。乙是 SOM。每个样本由最匹配的原型表示，同时利用网格邻域更新，尽量保留原数据的邻近关系。
\textbf{评分} 每种方法名称 0.5 分、产出的解释 0.5 分。把乙的更新对象写成仅最优原型，扣乙的解释分 0.5 分。
\end{enumerate}

\noindent\textbf{4. 线性自编码器、隐空间表示与 VAE}\pts{8}

\begin{enumerate}[label=(\arabic*)]
\item $z$ 为一维标量，$\widehat{\bm{x}}$ 为二维向量。两处维数各 $0.25$ 分。把瓶颈的维数当成重构维数，扣对应 $0.25$ 分。
\item
\begin{align*}
\|\bm{x}-\bm{u}\bm{u}^{\mathsf T}\bm{x}\|_2^2
&=\|\bm{x}\|_2^2-2(\bm{u}^{\mathsf T}\bm{x})^2
+\|\bm{u}\|_2^2(\bm{u}^{\mathsf T}\bm{x})^2\\
&=\|\bm{x}\|_2^2-(\bm{u}^{\mathsf T}\bm{x})^2.
\end{align*}
取期望，使用 $\E[\bm{x}\bm{x}^{\mathsf T}]=\bm{\Sigma}$，即得 $\operatorname{tr}(\bm{\Sigma})-\bm{u}^{\mathsf T}\bm{\Sigma}\bm{u}$。
展开和取期望各 $0.5$ 分。漏掉平方项中的交叉项，扣展开部分。
\item 写 $\bm{u}=a\bm{u}_1+b\bm{u}_2$，则 $a^2+b^2=1$ 且 $\bm{u}^{\mathsf T}\bm{\Sigma}\bm{u}=4a^2+b^2\le4$。极大值在 $\bm{u}=\pm\bm{u}_1$ 取得，最小损失为 $(4+1)-4=1$。线性瓶颈保留的是第一主成分所在的直线。
方向及 PCA 关系 $0.5$ 分，最小误差 $0.5$ 分。取小特征值方向，扣方向部分，并按后续计算步骤评分。
\item 使用加噪后的 $\widetilde{\bm{x}}$ 作为输入，原始的 $\bm{x}$ 作为目标。两项各 $0.25$ 分。把加噪数据作为目标扣相应 $0.25$ 分。
\item 稀疏自编码器约束隐藏层的激活值，使单次表示或平均激活更稀疏。堆栈自编码器将前一编码层的输出作为下一编码层的输入。各 $0.25$ 分。仅说“让权重越多为零越好”，扣稀疏部分。
\item 隐变量插值点为 $(3/4)(-2)+(1/4)6=0$，解码为 $g(0)=0$。若先解码再插值，则为 $(3/4)\cdot4+(1/4)\cdot36=12$。两个结果不同，因为解码函数是非线性的。
隐变量插值及解码 $0.5$ 分，端点解码后插值及比较 $0.5$ 分。只计算端点解码值而未比较，扣第二部分。
\item 两者目标为 $-3-\beta$ 与 $-5$，差值为 $2-\beta$。$0<\beta<2$ 时第一个更高，$\beta=2$ 时相同，$\beta>2$ 时第二个更高。
列出两个目标 $0.5$ 分，给对三个范围 $0.5$ 分。将重构对数似然当成越小越好的误差，扣比较部分。
\item 每个样本的先验 KL 都非负，平均值为零意味着验证集上的 $q(z\mid x)=p(z)$（几乎处处），编码器未通过该后验保留样本差异，这是后验坍塌的信号。要检查解码器，固定它的其他输入（若有），改变 $z$，比较输出分布或重构是否变化。
解释编码器的含义和给出检查方法各 $0.5$ 分。由零 KL 直接断言解码器必然恒定，扣检查部分。
\item 此时 $z=1$。忽略与参数无关的常数，$\ln p(x\mid z)=-\tfrac12(x-z^2)^2+\mathrm{const}$，因此
\[
\frac{\partial}{\partial z}\ln p(x\mid z)=2z(x-z^2)=4,
\qquad \frac{\partial}{\partial\mu}\ln p(x\mid z)=4,
\qquad \frac{\partial}{\partial s}\ln p(x\mid z)=2.
\]
$z$ 与对数似然函数 $0.5$ 分，对 $z$ 的链式求导 $0.5$ 分，对 $\mu,s$ 的梯度各 $0.25$ 分。遗漏 $\partial z/\partial s=\epsilon$，扣 $s$ 梯度的 $0.25$ 分。
\end{enumerate}

\noindent\textbf{5. 空洞卷积、门控网络与 Transformer 解码器}\pts{8}

\begin{enumerate}[label=(\arabic*)]
\item 等效跨度为 $(3-1)\cdot2+1=5$，每个方向实际采样三个位置。输出高 $7-5+1=3$，宽 $9-5+1=5$。跨度及输出各 $0.25$ 分。按普通 $3\times3$ 核计算，扣相应部分。
\item 每通道在自身的四个空间位置上取平均，得 $((2+0+4+2)/4,(1+3-1+9)/4)=(2,3)$。两个通道各 $0.25$ 分。跨所有通道求一个平均，扣 $0.5$ 分。
\item 拼接结果为 $(1,2,-1,3)$。后向状态由当前输入和后续时间的后向状态递推得到，通常依赖当前位置之后的输入。实时预测时这些输入尚不可用。
拼接和未来依赖各 $0.5$ 分。用逐元素相加代替拼接，扣拼接部分。
\item $r_t=\sigma(\ln3)=3/4$，$z_t=\sigma(0)=1/2$，所以
\[
\widetilde h_t=\tanh\!\bigl(\tfrac12\ln3-\tfrac34+\tfrac34\cdot1\bigr)=\tfrac12,
\qquad h_t=\tfrac12\cdot1+\tfrac12\cdot\tfrac12=\tfrac34.
\]
两门合计 $0.5$ 分，候选状态 $0.5$ 分，新状态 $0.5$ 分。漏用重置门，扣候选状态部分。误把候选状态当最终状态，扣新状态部分。
\item 因为 $h_t=o_t\odot\tanh(c_t)$，当前隐藏输出变成零。输出门不参与这一步的细胞状态更新，故细胞状态保持按原输入门、遗忘门和候选值算出的结果。
两种状态各 $0.25$ 分。把细胞状态也清零，扣细胞状态部分。
\item 训练依次输入 $\mathrm{BOS},u,v$。推理时第二次输入为模型刚预测出的 $w$。
训练移位的输入序列 $0.5$ 分，推理输入 $0.5$ 分。将 $\mathrm{EOS}$ 当作第三次输入，扣训练序列部分。
\item 本步上下文为 $(1/2)(2,0)+(1/3)(0,2)+(1/6)(2,2)=(4/3,1)$。下一步解码状态变化，注意力权重及上下文可以改变。
计算与关于下一步的解释各 $0.25$ 分。把所有编码状态等权平均，扣计算部分。
\item 查询由目标侧表示生成，键和值来自编码器的源侧表示。矩阵尺寸为
\[
Q:3\times2,\qquad K:5\times2,\qquad V:5\times3,
\qquad QK^{\mathsf T}:3\times5,\qquad \text{加权输出}:3\times3.
\]
来源 $0.5$ 分，尺寸 $0.5$ 分。互换源和目标会影响来源及相关尺寸，按对应部分扣分。
\item 令 $\operatorname{LN}$ 表示 LayerNorm，可写为
\begin{align*}
A&=\operatorname{LN}(X+\operatorname{MaskedAttn}(X)),\\
B&=\operatorname{LN}(A+\operatorname{CrossAttn}(A,H)),\\
C&=\operatorname{LN}(B+\operatorname{FFN}(B)).
\end{align*}
三子层顺序及交叉注意力来源 $0.5$ 分，三次残差和归一化位置 $0.5$ 分。把归一化移到残差相加前，扣后一部分。
\item 标准 FFN 独立地对每个位置的向量使用同一个前馈网络，其他位置的输入不出现在本位置的计算中，因而不会在这个子层直接传播。
结论与逐位置机制各 $0.25$ 分。仅以“各位置共享参数”推断跨位置通信，扣该问分。
\end{enumerate}

\noindent\textbf{6. 实验记录分类与判别面}\pts{8}
\begin{enumerate}[label=(\arabic*)]
\item “模拟”在 $A,B$ 中的概率是 $4/13,2/9$。后验赔率等于先验赔率乘似然比，即 $2(18/13)^2=648/169$。
\textbf{评分} 两个词概率 0.5 分，后验赔率 0.5 分。漏先验赔率扣 0.5 分。
\item 只出现一次时赔率为 $2(18/13)=36/13$。增加一次“模拟”会再乘因子 $18/13$，从 $36/13$ 变为 $648/169$。
\textbf{评分} 新赔率 0.5 分，重复词的倍乘解释 0.5 分。若把因子反过来写为 $13/18$，扣解释项 0.5 分。
\item 文档指示向量为 $(0,1,0)$。$A$ 下似然为 $(1-5/6)(3/6)(1-3/6)=1/24$，$B$ 下为 $(1-2/4)(2/4)(1-3/4)=1/16$。赔率为 $2(1/24)/(1/16)=4/3$。删去一次并不改变“模拟”存在的事件。
\textbf{评分} 两个似然 0.5 分，赔率与删词影响 0.5 分。忽略任何一个缺席词的补概率，扣似然项 0.5 分。
\item 未平滑时 $P(\text{丙}\mid A)=0$，相应后验为 0。加一后 $P(\text{丙}\mid A)=1/7,\ P(\text{丙}\mid B)=2/5$，故后验为 $(2/7)/(2/7+2/5)=5/12$。
\textbf{评分} 未平滑 0.5 分，平滑后验 0.5 分。平滑分母只加样本中实际出现的水平数，扣平滑项 0.5 分。
\item 用 $\mu_j=S_j/n$ 和 $\sigma_j^2=Q_j/n-\mu_j^2$，其中 $S_j,Q_j$ 分别为样本和、平方和。$A$ 的均值、方差为 $(0,2),(1,4)$，$B$ 为 $(2,-1),(1,4)$。
\textbf{评分} 均值 0.5 分，方差 0.5 分。把矩估计写成除以 $n-1$，扣方差项 0.5 分。
\item 以共享方差写出对数优势并消去二次项，
\[
 \ln 2-\tfrac12\!\left[
 x_1^2-(x_1-2)^2+\frac{(x_2-2)^2-(x_2+1)^2}{4}
 \right]
 =\ln 2+\tfrac{13}{8}-2x_1+\tfrac34 x_2.
\]
令右侧等于 0 即得直线。
\textbf{评分} 二次项抵消和对数优势 0.5 分，正确的等后验直线 0.5 分。若漏先验项，线性化过程正确可给第一项的方法分，判别式项扣 0.5 分。
\item 代入 $(1,0)$ 得后验对数优势 $\ln 2-3/8>0$，$P(A\mid x)=(1+e^{3/8}/2)^{-1}\approx0.5789$，判为 $A$。
\textbf{评分} 代入与分类 0.5 分，后验概率 0.5 分。直接把对数优势当概率扣 0.5 分。
\item $x_2^2$ 的系数是 $-1/(2\cdot4)+1/(2\cdot8)=-1/16$。$x_1$ 的一次项仍在，边界由含 $x_2^2$ 与 $x_1$ 的二次方程给出。
\textbf{评分} 系数 0.5 分，边界判断及理由 0.5 分。把方差再平方或符号写反，扣系数项 0.5 分。
\end{enumerate}

\noindent\textbf{7. 多项式回归的稀疏估计与多类逻辑回归}\pts{8}
\begin{enumerate}[label=(\arabic*)]
\item $\bm\varphi(x)=(1,x,x^2)^\top$，$w_2=0$ 时退化。遗漏截距特征或写错平方项扣 0.25 分，选错系数扣 0.25 分。
\item MLE 最大化 $p(\bm y\mid X,\bm w)$（0.5 分）。MAP 最大化 $p(\bm y\mid X,\bm w)p(w_1)p(w_2)$，即把先验作为额外权重（0.5 分）。写反似然条件方向扣前 0.5 分，MAP 漏先验或把先验直接与似然相加扣后 0.5 分。
\item 负对数后，与参数有关的项为
$\frac{1}{2\sigma^2}\sum_i(y_i-\hat y_i)^2+(|w_1|+|w_2|)/s$（0.5 分）。
两边乘以正数 $\sigma^2$，所以 $\lambda=\sigma^2/s$（0.5 分）。误对截距加惩罚或写错负对数项扣前 0.5 分，系数倒置扣后 0.5 分。
\item 对每个坐标，分别最小化 $\frac12(w-a)^2+|w|$。$w>0$ 时令 $w-a+1=0$，$w<0$ 时令 $w-a-1=0$。$w=0$ 的次梯度区间为 $[-a-1,-a+1]$。$a=2$ 得 $w_1=1$，$a=0.5$ 时零点满足条件，得 $w_2=0$。两坐标各 0.5 分。该坐标只给结果而无对应半轴或零点检验，扣 0.25 分。结果错扣 0.5 分。
\item 可行域是顶点 $(\pm1,0),(0,\pm1)$ 的菱形（0.5 分）。损失等高线可能首先碰到轴上的尖角，在这些位置另一坐标为零（0.5 分）。画成圆形扣前 0.5 分，声称每个 Lasso 解都有零分量或缺少尖角解释扣后 0.5 分。
\item 指数值是 $(3,2,1)$，归一化得 $\bm p=(1/2,1/3,1/6)^\top$（0.5 分）。$L=-\ln p_2=\ln 3\approx1.099$（0.5 分）。概率归一化错误扣前 0.5 分，使用错误标签算损失扣后 0.5 分。
\item $\partial(\ln\sum_j e^{a_j})/\partial a_k=e^{a_k}/\sum_j e^{a_j}=p_k$（0.5 分）。减去第二项的导数，得 $\partial L/\partial a_k=p_k-\bm 1_{\{k=2\}}$（0.5 分）。漏写指数项导数或漏标签项，各扣对应的 0.5 分。
\item $W$ 的第二行梯度为 $(p_2-1)\bm x^\top=(-4/3,2/3)$。缺少输入因子扣 0.25 分，数值或正负号错误扣 0.25 分。
\item 不改变。分子分母都乘 $e^c$，约去即可。判断和理由各 0.25 分。
\item $\ln(p_1/p_2)=a_1-a_2=(W_{1,:}-W_{2,:})\bm x+b_1-b_2$（0.25 分）。题设点取值 $\ln(3/2)\approx0.4055$（0.25 分）。方向或数值错误，扣对应分数。
\end{enumerate}

\noindent\textbf{8. 浮标观测的缩放前向计算}\pts{8}
\begin{enumerate}[label=(\arabic*)]
\item $\widetilde\alpha_1(i)=\pi_i b_i(o_1)$。对 $t>1$，$\widetilde\alpha_t(j)=b_j(o_t)\sum_i\widehat\alpha_{t-1}(i)a_{ij}$。$c_t=\sum_j\widetilde\alpha_t(j)$，$\widehat\alpha_t(j)=\widetilde\alpha_t(j)/c_t$。
\textbf{评分} 初始化及递推 0.5 分，正确归一化 0.5 分。将前一时刻的未归一化向量用于本题定义的递推，扣递推项 0.5 分。
\item $\widetilde\alpha_1=(2,2,3)/16$，所以 $c_1=7/16$，$\widehat\alpha_1=(2,2,3)/7$。
\textbf{评分} 两项各 0.5 分。只做转移或漏乘发射导致的错误，扣对应项的分。
\item $\widetilde\alpha_2=(20,27,9)/112$，所以 $c_2=1/2$，$\widehat\alpha_2=(20,27,9)/56$。
\textbf{评分} 两项各 0.5 分。后一向量若按本人正确的计算方法归一化，可给对应方法分。
\item $\widetilde\alpha_3=(130,76,249)/896$，所以 $c_3=65/128$，$\widehat\alpha_3=(130,76,249)/455$。
\textbf{评分} 两项各 0.5 分。分量和不为 1 的向量不能得归一化向量项的分。
\item $\widetilde\alpha_4=(1408,585,1593)/7280$，所以 $c_4=1793/3640$，$\widehat\alpha_4=(1408,585,1593)/3586$。
\textbf{评分} 两项各 0.5 分。仍使用 $Y$ 的发射概率，在相关项扣分。
\item $P(O\mid\lambda)=\prod_t c_t=1793/32768\approx0.05472$，$\ln P(O\mid\lambda)=\sum_t\ln c_t\approx-2.906$。
\textbf{评分} 乘积及概率 0.25 分，对数和及其数值 0.25 分。照搬未缩放方法对 $\widehat\alpha_4$ 直接求和，扣概率项 0.25 分。
\item 最后一个观测产生的后向贡献为 $(7,9,8)/16$，除以 $c_4$ 得
\[
 \widehat\beta_3=(3185,4095,3640)/3586.
\]
在此缩放约定下，状态后验为
\[
 \gamma_3(i)=\widehat\alpha_3(i)\widehat\beta_3(i),\qquad
 P(S_3=3\mid O,\lambda)=996/1793.
\]
\textbf{评分} 缩放后向量 0.5 分，平滑后验 0.5 分。漏除 $c_4$ 扣后向量项 0.5 分。
\item 下一次的预测概率为 $\sum_j b_j(X)\sum_i\widehat\alpha_4(i)a_{ij}=27865/57376\approx0.4857$。
\textbf{评分} 正确进行转移再发射 0.25 分，计算结果 0.25 分。跳过转移直接乘发射扣第一项。
\item 长序列的联合概率可能很小，直接连续相乘会下溢。逐步归一化使当前向量保持可计算的尺度，保存 $\ln c_t$ 的和即可恢复对数似然。离散 HMM 的隐藏状态取有限类别，由随机状态转移和发射分布建模。普通 RNN 的隐藏状态通常是连续向量，由前一向量和当前输入经确定的递推映射计算。
\textbf{评分} 缩放与对数的作用 0.5 分，HMM/RNN 对隐藏状态的区别 0.5 分。只说缩放让最后的原始概率变大，不给缩放项分。
\end{enumerate}
